\section{总结与展望/ConclusIon And RecommendAtIons}

随着无线频谱资源的日益紧缺以及B5G/6G网络对全场景环境感知的迫切需求，集成感知与通信（ISAC）技术已成为突破传统通信范式的关键基础设施。ISAC通过打破雷达探测与数据通信之间的技术边界，在统一的硬件架构、频谱资源和信号处理框架下实现双重功能，有效解决了硬件冗余、频谱碎片化以及资源竞争等痛点，是实现未来智能世界“万物智联”与“深度感知”融合的关键技术路径。

本文围绕ISAC技术的核心要素，从五个维度系统性地探讨了其关键技术体系。首先，是演进历程与设计愿景，介绍了ISAC从雷达通信共存向深度融合演进的基本概念及其在6G愿景中的核心地位。接着在理论基石层面，探讨了通感一体化的信息论极限与性能权衡，通过分析检测概率、均方误差（MSE）与通信容量等指标，揭示了如何在有限资源下实现感知精度与传输速率的最优帕累托边界。其次，在信号与波形层面，深入分析了基于OFDM、SSB及PRS等现有波形的改造与联合波形设计，探讨了匹配滤波增益、自干扰抑制及未知数据负荷下的信号提取算法，强调了波形统一性对系统集成度的贡献。再次，在架构与网络层面，探讨了感知网络（Perceptive Networks）的构建，重点介绍了帧级ISAC与网络级分布式感知的协同机制，以及如何利用云无线接入网实现大尺度的环境重构。最后，在应用场景与验证层面，阐述了ISAC在车联网（V2X）协同编队、高精度同步定位与建图（SLAM）及人类活动识别等领域的落地潜力，为构建全感知无线网络提供了从理论到实践的完整闭环。

%展望未来，ISAC技术的大规模商用仍需在以下几个方向持续探索。一是高频段与超大规模MIMO的融合，利用毫米波及太赫兹（THz）极大的带宽提升感知分辨率，并借助空间自由度克服复杂的非视距传播。二是人工智能赋能的语义感知，将深度学习引入感知信号处理，提升复杂多径环境下目标特征提取的鲁棒性，实现从“信号探测”向“语义理解”的跨越。三是通感融合的标准化与安全性，建立统一的空口协议标准以支持跨厂商设备的协同，并研究感知数据的隐私保护机制，确保在获取环境信息的同时不侵犯用户安全。
